\section{Startup Minimum Viable Safety}

Las grandes plataformas pueden formar equipos dedicados de trust-and-safety. Las empresas emergentes, en cambio, suelen enfrentar riesgos desde que aparece el contenido generado por usuarios, pero les faltan experiencia, presupuesto y prioridad para atenderlos. Esperar a «gobernar cuando haya más escala» convierte en deuda técnica el modelo de datos, los registros, el flujo de reportes y la cultura. Esta sección presenta la minimum viable safety: no se trata de copiar la organización de una gran empresa, sino de establecer un ciclo cerrado mínimo que pueda ampliarse.

\subsection{Risk inventory, known-content control e incident readiness}

Primero, el equipo de una empresa emergente debe enumerar varios puntos: quién puede interactuar con quién dentro del producto, qué contenido puede subirse, si los menores de edad pueden verlo, cómo se propagan los mensajes privados y los grupos, y qué modelos o servicios de almacenamiento de terceros participan en el procesamiento. Después, elige la combinación mínima de coincidencia con contenido conocido, reportes de usuarios, moderation queue, account action, audit log y emergency escalation. También define de forma explícita qué se resolverá con build, con buy o con specialist partnership.

\lecturefigure{15-startup-maturity.png}{Las empresas emergentes necesitan pasar del risk inventory a una ruta de minimum viable safety que pueda ampliarse.}{Subtítulos locales 35:20--37:20; redibujo conceptual basado en la discusión de la clase sobre la startup barrier.}

\begin{knowledgebox}{Lectura de la figura: primero, que la responsabilidad y los incidentes tengan a dónde ir}
El risk inventory responde cuál es la superficie de exposición del producto. Los minimum controls aportan known-hash, report y basic review. El managed program agrega predictive triage, policy, metrics y auditoría. Por último, la proactive prevention incorpora threat research, red teaming y protección de usuarios. En la primera etapa, lo más importante no es comprar todas las herramientas. Lo que importa es designar un owner, conservar los eventos necesarios, establecer una vía de reporte externo y, ante un incident real, poder localizar a la persona de guardia y el registro de decisiones.
\end{knowledgebox}

\begin{importantbox}{Minimum viable safety checklist}
\begin{enumerate}
  \item Un safety owner con poder de decisión y una on-call escalation;
  \item el risk inventory del producto, el modelo de edad e interacción y la política de uso aceptable;
  \item una ruta de report/block/help visible y fácil de usar para los usuarios;
  \item known-content control, basic queue, case log y deduplicación de eventos repetidos;
  \item minimización de datos, permisos de acceso, retention e incident playbook;
  \item SLA de proveedores, degradación ante fallas, reportes legalmente obligatorios y lista de contactos de instituciones especializadas.
\end{enumerate}
\end{importantbox}

\begin{warningbox}{«Hacerlo después» amplifica el costo de migración}
Si desde el inicio no hay ID estables de user/account/content, más adelante no es posible agregar los eventos. Sin audit log, no se pueden explicar las medidas que se tomaron antes. Sin un age-aware product model, la única opción es imponer restricciones a la fuerza una vez terminada la UI. Sin supplier contract, es difícil rastrear los datos sensibles. Cuanto más tarde se incorpore la safety architecture, más probable es que obligue a reescribir identity, storage, messaging y el moderation core.
\end{warningbox}

\teachervoice{Quien expone resume los obstáculos de las startups en tres factores: awareness, money y priority. Con esto evita limitarse a culpar a los fundadores. Para la ingeniería, la clase concluye que hay que ofrecer infraestructura compartida, una ruta de madurez clara y servicios asequibles. Así, los equipos pueden empezar con un ciclo cerrado mínimo y no tienen que elegir entre «no hacer nada» y «construir por su cuenta un departamento de seguridad completo».}

\subsection{Resumen de la sección}

La seguridad en una empresa emergente empieza por tener claros los riesgos del producto, los responsables y las rutas de reporte e incidentes. Después se incorporan de forma gradual la detección, la evaluación y la prevención. Integrar esto desde el principio suele costar menos que limpiar después de un incidente, y también facilita conservar los límites correctos de datos y permisos.
